\documentclass[11pt]{article}
\usepackage[margin=1in]{geometry}
\usepackage{amsmath,amssymb,amsthm,mathtools,bm}
\usepackage{booktabs,longtable,array}
\usepackage[dvipsnames]{xcolor}
\usepackage[colorlinks=true,linkcolor=MidnightBlue,citecolor=BrickRed,urlcolor=MidnightBlue]{hyperref}
\usepackage{microtype}

\hypersetup{
  pdftitle={Explicit SU(2) Heat-Kernel Initialization},
  pdfauthor={},
  pdfsubject={One-plaquette derivative and exceptional-event certificates},
  pdfkeywords={Yang--Mills, SU(2), heat kernel, Poisson summation, large fields}
}

\newtheorem{theorem}{Theorem}[section]
\newtheorem{proposition}[theorem]{Proposition}
\newtheorem{lemma}[theorem]{Lemma}
\newtheorem{corollary}[theorem]{Corollary}
\theoremstyle{definition}
\newtheorem{definition}[theorem]{Definition}
\newtheorem{remark}[theorem]{Remark}

\newcommand{\R}{\mathbb R}
\newcommand{\Id}{\mathrm I}
\newcommand{\SU}{\mathrm{SU}(2)}
\newcommand{\statusproved}{\textcolor{ForestGreen}{\textsc{proved}}}
\newcommand{\statusnumeric}{\textcolor{RoyalBlue}{\textsc{computed}}}
\newcommand{\statusopen}{\textcolor{BrickRed}{\textsc{open}}}

\title{Explicit \texorpdfstring{$SU(2)$}{SU(2)} Heat-Kernel Initialization\\
\large Local derivative budgets, image-tail control, and a gauge-invariant exceptional event}
\author{Research dossier---one-plaquette theorem and block-interface audit}
\date{20 July 2026}

\begin{document}
\maketitle

\begin{abstract}
The preceding one-step effective-action theorem propagates five declared
derivative budgets but does not initialize them for the exact nonlinear gauge
action.  This paper removes the existential constant from the radial
one-plaquette part of that problem.  With $SU(2)$ identified with the unit
three-sphere and heat generator $\frac12\Delta$, the character and
Poisson-image representations are matched exactly.  Pairing positive and
negative images yields explicit uniform bounds on the first four derivatives
of the image correction, including an analytic omitted-image tail.  For
$0<t\le0.2$ and class angle $\theta\le1$, the resulting logarithmic-curvature
certificate is
\[
 t\left\|\nabla^2\left(-\log p_t-\frac{\theta^2}{2t}\right)\right\|
 \le0.0824566,
 \qquad
 t\left\|D^3\left(-\log p_t-\frac{\theta^2}{2t}\right)\right\|
 \le0.149094.
\]
The event $\theta(U_p)\ge\rho$ is conjugation invariant.  Its one-plaquette
probability is enclosed by an analytically integrated character series with a
geometric truncation bound and explicit floating-roundoff allowance.  A
\statusproved\ chain-rule theorem transports the radial constants through any
curvature chart whose first three derivative norms are supplied.  Establishing
those norms on the full $2^4$ quotient chart and converting plaquette tails
into dependent block-polymer activities remain \statusopen.
\end{abstract}

\tableofcontents

\section{Target, normalization, and scope}

The one-step theorem of the preceding stage is exact but typed: it accepts
numbers such as $\varepsilon_F,\varepsilon_X,\varepsilon_Y,\tau_F,\tau_X$ and
returns a generated-action certificate.  The earliest missing implication is
to obtain such numbers from the nonlinear gauge action.  This stage isolates
the first factor in that implication:
\[
 \text{exact radial heat-kernel action}
 \longrightarrow
 \text{curvature-coordinate derivative budgets}.
\]
The heat-kernel lattice action follows Menotti and Onofri
\cite{MenottiOnofri}.  Explicit heat-kernel eigenfunction and image
representations on homogeneous spaces are part of the harmonic-analysis
framework developed by Camporesi \cite{Camporesi}.

Identify $SU(2)$ with the unit sphere $S^3$.  Every conjugacy class has a
representative
\[
 U(\theta)=\operatorname{diag}(e^{i\theta},e^{-i\theta}),
 \qquad 0\le\theta\le\pi.
\]
Haar probability has radial density
\begin{equation}\label{eq:haar}
 d\mu_{\mathrm{rad}}(\theta)=\frac{2}{\pi}\sin^2\theta\,d\theta.
\end{equation}
We fix heat generator $\frac12\Delta_{S^3}$.  The eigenvalue in the
$(n+1)$-dimensional irreducible representation is $n(n+2)$, so the central
heat kernel is
\begin{equation}\label{eq:character}
 p_t(\theta)=\sum_{n=0}^{\infty}(n+1)e^{-n(n+2)t/2}
 \frac{\sin((n+1)\theta)}{\sin\theta}.
\end{equation}
This convention fixes the principal local action as $\theta^2/(2t)$.  Changing
the metric or Laplacian convention changes every numerical time value and must
be performed before comparison with another regulator.

The target is one plaquette and one logarithmic curvature chart strictly
inside the cut locus.  The result is C2: a rigorous restricted theorem.  It is
not a full block theorem, a large-field polymer theorem, or a mass-generation
result.  The Clay target remains the continuum construction and Hamiltonian
gap described by Jaffe and Witten \cite{JaffeWitten}.

\section{Poisson-image representation}

Set $m=n+1$ in \eqref{eq:character}.  Since $n(n+2)=m^2-1$,
\[
 p_t(\theta)=\frac{e^{t/2}}{\sin\theta}
 \sum_{m=1}^{\infty}m e^{-m^2t/2}\sin(m\theta).
\]
Poisson summation of the periodic Gaussian, followed by differentiation in
$\theta$, gives
\begin{equation}\label{eq:image}
 p_t(\theta)=C(t)\frac{1}{\sin\theta}
 \sum_{k\in\mathbb Z}(\theta+2\pi k)
 \exp\left[-\frac{(\theta+2\pi k)^2}{2t}\right],
 \qquad
 C(t)=\frac{e^{t/2}\sqrt{2\pi}}{2t^{3/2}}.
\end{equation}

Factor the principal image:
\begin{equation}\label{eq:q-def}
 \sum_{k\in\mathbb Z}(\theta+2\pi k)e^{-(\theta+2\pi k)^2/(2t)}
 =\theta e^{-\theta^2/(2t)}(1+q_t(\theta)).
\end{equation}
Pairing $k$ with $-k$, and writing $a_k=2\pi k$, yields the analytic even
function
\begin{equation}\label{eq:q-pair}
 q_t(\theta)=2\sum_{k=1}^{\infty}e^{-a_k^2/(2t)}
 \left[
 \cosh\left(\frac{a_k\theta}{t}\right)
 -\frac{a_k}{\theta}\sinh\left(\frac{a_k\theta}{t}\right)
 \right].
\end{equation}
The apparent singularity at $\theta=0$ is removable.  Indeed,
\begin{equation}\label{eq:sinh-integral}
 \frac{a_k}{\theta}\sinh\left(\frac{a_k\theta}{t}\right)
 =\frac{a_k^2}{t}\int_0^1
 \cosh\left(\frac{s a_k\theta}{t}\right)ds.
\end{equation}

\section{Explicit image derivative bounds}

Fix $0<\rho<\pi/2$ and $0<t\le t_0$.  For $j=0,1,2,3$, define
\begin{equation}\label{eq:Qj}
 Q_j(t_0,\rho):=
 2\sum_{k=1}^{\infty}
 e^{-a_k^2/(2t_0)+a_k\rho/t_0}
 \left(\frac{a_k}{t_0}\right)^j
 \left(1+\frac{a_k^2}{t_0(j+1)}\right).
\end{equation}

\begin{lemma}[Uniform image bound, \statusproved]\label{lem:image-bound}
Let
\begin{equation}\label{eq:time-condition}
 t_0\le \frac{(2\pi)^2-4\pi\rho}{8}.
\end{equation}
Then
\[
 \sup_{0<t\le t_0}\sup_{0\le\theta\le\rho}
 |q_t^{(j)}(\theta)|\le Q_j(t_0,\rho),
 \qquad j=0,1,2,3.
\]
After retaining images $1\le k\le K$, the omitted series is bounded by the
first omitted term divided by $1-r_{j,K}$, where
\begin{equation}\label{eq:image-ratio}
 r_{j,K}=\left(\frac{K+2}{K+1}\right)^{j+2}
 \exp\left[-\frac{(2\pi)^2(2K+3)}{2t_0}
 +\frac{2\pi\rho}{t_0}\right]<1.
\end{equation}
\end{lemma}
\begin{proof}
Differentiate \eqref{eq:q-pair} and use
\eqref{eq:sinh-integral}.  The $j$th derivative introduces
$(a_k/t)^j$ and $s^j$ under the integral.  Bounding hyperbolic sine and cosine
by $e^{a_k\rho/t}$ and integrating $s^j$ gives the $k$th summand of
\eqref{eq:Qj}.  Each summand has the form $t^{-p}e^{-c_k/t}$ with
$p\le j+1\le4$.  Condition \eqref{eq:time-condition}, using the first image,
makes it increasing on $(0,t_0]$, so $t=t_0$ is the uniform endpoint.  The
polynomial ratio between consecutive $k$ terms is at most
$((k+1)/k)^{j+2}$; the Gaussian ratio gives \eqref{eq:image-ratio}, which
decreases with $k$.  Geometric summation proves the tail bound.
\end{proof}

Write
\begin{equation}\label{eq:action-factor}
 s_t(\theta):=-\log p_t(\theta)
 =\frac{\theta^2}{2t}+a(\theta)+b_t(\theta)+c(t),
\end{equation}
where
\[
 a(\theta)=\log\frac{\sin\theta}{\theta},
 \qquad b_t(\theta)=-\log(1+q_t(\theta)).
\]
If $Q_0<1$, direct differentiation gives
\begin{align}
 B_2&:=\sup|b_t''|
 \le \frac{Q_2}{1-Q_0}+\frac{Q_1^2}{(1-Q_0)^2},\label{eq:B2}\\
 B_3&:=\sup|b_t'''|
 \le \frac{Q_3}{1-Q_0}
 +\frac{3Q_1Q_2}{(1-Q_0)^2}
 +\frac{2Q_1^3}{(1-Q_0)^3}.\label{eq:B3}
\end{align}

The geometric term has
\begin{equation}\label{eq:A23}
 A_2(\rho):=\sup_{[0,\rho]}|a''|
 =\csc^2\rho-\rho^{-2},
 \qquad
 A_3(\rho):=\sup_{[0,\rho]}|a'''|
 =2\rho^{-3}-\frac{2\cos\rho}{\sin^3\rho}.
\end{equation}
The equalities follow from the positive-coefficient expansions of
$\csc^2\theta-\theta^{-2}$ and its derivative on this interval.

\section{Local curvature-action certificate}

Let $X\in\mathfrak{su}(2)\cong\R^3$ and $\theta=\|X\|$.  A radial function
$u(\theta)$ has Hessian eigenvalues $u''(\theta)$ in the radial direction and
$u'(\theta)/\theta$ in the two tangential directions.

\begin{theorem}[Explicit local heat-kernel initializer, \statusproved]
\label{thm:initializer}
Assume Lemma~\ref{lem:image-bound}, $Q_0<1$, and $\|X\|\le\rho$.  Then for all
$0<t\le t_0$,
\begin{equation}\label{eq:epsilon-hk}
 t\left\|\nabla_X^2s_t(X)-\Id/t\right\|
 \le \varepsilon_{\mathrm{HK}}
 :=t_0(A_2+B_2).
\end{equation}
Moreover,
\begin{equation}\label{eq:gradient-hk}
 t\left\|\nabla_X\left(s_t(X)-\frac{\|X\|^2}{2t}\right)\right\|
 \le g_{\mathrm{HK}}:=\rho\varepsilon_{\mathrm{HK}},
\end{equation}
and
\begin{equation}\label{eq:tau-hk}
 t\left\|D_X^3\left(s_t(X)-\frac{\|X\|^2}{2t}\right)\right\|
 \le\tau_{\mathrm{HK}}:=4t_0(A_3+B_3).
\end{equation}
\end{theorem}
\begin{proof}
Apply \eqref{eq:B2} and \eqref{eq:A23} to the radial Hessian.  Since the
remainder derivative vanishes at zero, the mean-value theorem bounds the
tangential eigenvalue and gives \eqref{eq:gradient-hk}.  For a radial $C^3$
function $u$, set
\[
 \beta=\frac{u''}{\theta}-\frac{u'}{\theta^2}.
\]
If $M_3=\sup|u'''|$, integration gives $|\beta|\le M_3/2$.  The Euclidean
third derivative tensor is the sum of a radial tensor with coefficient
$u'''-3\beta$ and three mixed tensors with coefficient $\beta$.  Its operator
norm is therefore at most $|u'''-3\beta|+3|\beta|\le4M_3$, proving
\eqref{eq:tau-hk}.
\end{proof}

For $(t_0,\rho)=(0.2,1)$ the analytic quantities are
\begin{center}
\begin{tabular}{@{}lr@{}}
\toprule
Quantity & Value\\
\midrule
$A_2$ & $0.4122829274$\\
$A_3$ & $0.1863671214$\\
$Q_0$ & $2.3943\times10^{-27}$\\
$Q_3$ & $1.8840\times10^{-23}$\\
$\varepsilon_{\mathrm{HK}}$ & $0.08245658549$\\
$\tau_{\mathrm{HK}}$ & $0.1490936971$\\
\bottomrule
\end{tabular}
\end{center}
The image contribution is much smaller than the group-Jacobian correction in
this window; the implementation nevertheless retains its analytic tail.

\section{Composition with a curvature chart}

Let $X(A)$ be a declared plaquette-log map on a link-coordinate chart.  Write
$D_0=DX(0)$ and suppose
\begin{equation}\label{eq:chart-norms}
 \|D_0\|\le M_0,\quad \|DX\|\le M_1,\quad
 \|DX-D_0\|\le\delta_1,\quad
 \|D^2X\|\le M_2,\quad \|D^3X\|\le M_3.
\end{equation}

\begin{proposition}[Chart propagation, \statusproved]\label{prop:chart}
On a region where $\|X(A)\|\le\rho$, the Hessian remainder of
$s_t\circ X$ relative to $D_0^*D_0/t$ is bounded by
\begin{equation}\label{eq:chart-epsilon}
 \varepsilon_{\mathrm{chart}}
 =\delta_1(M_1+M_0)
 +\varepsilon_{\mathrm{HK}}M_1^2
 +(\rho+g_{\mathrm{HK}})M_2.
\end{equation}
Its scaled third derivative is bounded by
\begin{equation}\label{eq:chart-tau}
 \tau_{\mathrm{chart}}
 =\tau_{\mathrm{HK}}M_1^3
 +3(1+\varepsilon_{\mathrm{HK}})M_1M_2
 +(\rho+g_{\mathrm{HK}})M_3.
\end{equation}
\end{proposition}
\begin{proof}
Use
\[
 D^2(s_t\circ X)=DX^*(D^2s_t)DX+\langle Ds_t,D^2X\rangle
\]
and
$\|DX^*DX-D_0^*D_0\|\le\delta_1(M_1+M_0)$ to obtain
\eqref{eq:chart-epsilon}.  The third-order chain rule is
\[
 D^3(s_t\circ X)=D^3s_t[DX]^3
 +3D^2s_t[D^2X,DX]+Ds_t[D^3X].
\]
Theorem~\ref{thm:initializer} gives
$t\|D^2s_t\|\le1+\varepsilon_{\mathrm{HK}}$ and
$t\|Ds_t\|\le\rho+g_{\mathrm{HK}}$, proving
\eqref{eq:chart-tau}.
\end{proof}

For the identity curvature chart, these formulas return exactly
$\varepsilon_{\mathrm{HK}}$ and $\tau_{\mathrm{HK}}$.  For a block, a declared
aggregation operator norm must multiply the per-plaquette bounds.  It cannot
be silently replaced by either independence or raw plaquette count.  A stress
example with $(M_0,M_1,\delta_1,M_2,M_3)=(1,1.1,0.1,0.2,0.3)$ and aggregation
factor $6$ gives Hessian budget $3.16>2$.  Those chart values are illustrative,
not certified BCH constants; the failure demonstrates why their sharp
derivation is the next proof obligation.

\section{Gauge-invariant exceptional probability}

Define the one-plaquette exceptional event
\begin{equation}\label{eq:event}
 E_{\rho,p}=\{U_p:\theta(U_p)\ge\rho\}.
\end{equation}
Since class angle is invariant under $U_p\mapsto gU_pg^{-1}$, this event is
gauge invariant.

Combining \eqref{eq:character} and \eqref{eq:haar} gives
\begin{equation}\label{eq:probability}
 \mathbb P_t(E_{\rho,p})
 =\frac{2}{\pi}\sum_{n=0}^{\infty}(n+1)e^{-n(n+2)t/2}I_{n+1}(\rho),
\end{equation}
where
\begin{align}
 I_1(\rho)&=\frac{\pi-\rho}{2}+\frac{\sin2\rho}{4},\label{eq:I1}\\
 I_m(\rho)&=-\frac12\left(
 \frac{\sin((m-1)\rho)}{m-1}
 -\frac{\sin((m+1)\rho)}{m+1}
 \right),\qquad m\ne1.\label{eq:Im}
\end{align}

\begin{proposition}[Character-tail enclosure, \statusproved]
\label{prop:probability}
After summing $0\le n\le N$, the absolute omitted tail in
\eqref{eq:probability} is at most
\begin{equation}\label{eq:char-tail}
 R_N(t)=\frac{2(N+2)e^{-(N+1)(N+3)t/2}}{1-r_N(t)},
 \qquad
 r_N(t)=\frac{N+3}{N+2}e^{-(2N+5)t/2},
\end{equation}
provided $r_N(t)<1$.
\end{proposition}
\begin{proof}
Equations \eqref{eq:I1}--\eqref{eq:Im} are the elementary integrals of
$\sin(m\theta)\sin\theta$ over $[\rho,\pi]$.  Their absolute values are at most
$\pi$, so the omitted $n$th contribution is at most
$2(n+1)e^{-n(n+2)t/2}$.  The ratio of consecutive majorants decreases with
$n$ and equals $r_N(t)$ at the first omitted index.  Sum the geometric
majorant.
\end{proof}

At $\rho=1$ and $N=80$, the analytic truncation is far below the explicitly
recorded floating-roundoff allowance:
\begin{center}
\begin{tabular}{@{}lrr@{}}
\toprule
$t$ & $\mathbb P_t(\theta\ge1)$ & analytic series tail\\
\midrule
$0.20$ & $0.1474007238$ & $1.74\times10^{-290}$\\
$0.10$ & $0.01580367907$ & $1.69\times10^{-144}$\\
$0.05$ & $0.0001437428033$ & $1.70\times10^{-71}$\\
\bottomrule
\end{tabular}
\end{center}
The code reports an interval rather than the midpoint alone.  The probability
is a genuine one-plaquette large-field datum.  Neighboring plaquettes share
links; therefore these numbers cannot be tensorized into a block or polymer
activity without a dependence theorem.

\section{Executable verification}

The implementation performs fourteen positive and hostile checks:
\begin{itemize}
\item character and Poisson-image ratios agree to floating precision at three
nonzero class angles;
\item omitted image tails remain positive even below the ordinary floating
range, with their logarithmic bounds retained;
\item decreasing $t_0$ improves both local derivative budgets;
\item the character probability interval contains a longer-series reference;
\item probability is normalized at $\rho=0$ and decreases with the calibrated
radius;
\item the identity chart preserves the radial budgets exactly;
\item cut-locus radii, excessive time ranges, invalid chart norms, and
nonpositive aggregation factors fail closed.
\end{itemize}

This is executable analytic bookkeeping rather than formal proof-assistant
verification.  Image and character truncation errors are bounded analytically;
the partial sums use double precision with a separately recorded roundoff
allowance.

\section{Frontier audit and next theorem}

\begin{longtable}{@{}p{0.08\textwidth}p{0.65\textwidth}p{0.14\textwidth}@{}}
\toprule
ID & Required implication & Status\\
\midrule
\endhead
H1 & Exact radial heat-kernel derivative initializer and one-plaquette exceptional probability & \statusproved\\
F1 & Explicit transported plaquette-log, Coulomb-slice, retained/fiber projection, and overlap constants on the full $2^4$ block & \statusopen\\
F2 & Sharp block incidence aggregation preserving the Stage-06 fiber margin & \statusopen\\
F3 & Dependent gauge-invariant bad-block polymer estimate, including center flux, holonomy, exceptional orbits, boundaries, and topology & \statusopen\\
F4 & Initialization of all five Stage-06 components for generated localized interactions & \statusopen\\
F5 & Cutoff-uniform repeated RG closure and retained-base massive crossover & \statusopen\\
F6 & Nontrivial regulator-independent OS continuum theory and Hamiltonian mass gap & \statusopen\\
\bottomrule
\end{longtable}

The result changes the location of the first missing implication.  The radial
heat kernel itself no longer contributes an existential $C_B$: its geometric
and image corrections are explicit.  The main local uncertainty is now the
representation map from 216 noncommuting plaquette holonomies, through a gauge
slice and coarse retained data, into the $51+357$ quotient/fiber coordinates.
The next bounded theorem should derive those block constants using the exact
$SU(2)$ quaternion product and logarithm, rather than a formal BCH remainder.
It should then insert the resulting numbers into the Stage-06 inequalities.

Even if that local block closes, the one-plaquette tail must still be replaced
by a dependent polymer activity.  Large fields, topology, repeated scales,
infrared crossover, continuum existence, reflection positivity, observable
density, and spectral transfer remain indispensable before any mass-gap claim.

\begin{thebibliography}{9}
\bibitem{MenottiOnofri}
P.~Menotti and E.~Onofri,
\emph{The action of $SU(N)$ lattice gauge theory in terms of the heat kernel on
the group manifold},
Nucl. Phys. B \textbf{190} (1981), 288--300,
\url{https://doi.org/10.1016/0550-3213(81)90560-5}.

\bibitem{Camporesi}
R.~Camporesi,
\emph{Harmonic analysis and propagators on homogeneous spaces},
Phys. Rep. \textbf{196} (1990), 1--134,
\url{https://doi.org/10.1016/0370-1573(90)90120-Q}.

\bibitem{Balaban1987}
T.~Balaban,
\emph{Renormalization group approach to lattice gauge field theories: I.
Generation of effective actions in a small field approximation and a coupling
constant renormalization in four dimensions},
Commun. Math. Phys. \textbf{109} (1987), 249--301,
\url{https://doi.org/10.1007/BF01215223}.

\bibitem{JaffeWitten}
A.~Jaffe and E.~Witten,
\emph{Quantum Yang--Mills Theory}, Clay Mathematics Institute Millennium
Problem description,
\url{https://www.claymath.org/wp-content/uploads/2022/06/yangmills.pdf}.
\end{thebibliography}

\end{document}
